
{\Large TEST} \\
\begin{figure}[h!]
\centering
\includegraphics[width=\textwidth]{chaos}
\caption{test}
\label{fig:test1}
\end{figure}

%\begin{lstlisting}
\lstinputlisting[language=C++, firstline=9, lastline=18, caption=Painter.cpp;9:18, label=Painter.cpp]{Painter.cpp}
%\lstlistingname{Painter.cpp;9:18}
%\end{lstlisting}

\newpage

\section{Wstęp}

\subsection{Motywacja}
\paragraph{}
Symulacja płynów jest stosunkowo złożonym i skomplikowanym problemem. W ostatnim dziesięcioleciu staje się jednak coraz bardziej popularnym tematem opracowań naukowych. Obliczeniowa mechanika płynów (ang. \textbf{C}omputational \textbf{F}luid \textbf{D}ynamics) wykorzystuje metody numeryczne oraz stale rosnącą moc obliczeniową komputerów do rozwijania tej dyscypliny. Powodem złożoności zachowania płynów jest skomplikowany mechanizm działania różnych zjawisk składających się na pełen <<obraz>> płynów, takich jak turbulentność, laminarność czy napięcie powierzchniowe płynu.\\
Symulacje płynu można przeprowadzać w dwojaki sposób: offline i online.
Metody offline zapewniają największą dokładność obliczeń. Oprogramowania stosujące te techniki wykorzystywane są m.in. w przemyśle motoryzacyjnym, lotniczym i kosmicznym gdzie precyzja kalkulacji jest bardzo istotna. Metody dokładne sprowadzają się do rozwiązania równań Naviera-Stokesa.\\
W ramach tej pracy wykonano symulacje zachowania płynów w oparciu o metodę online -- SPH (Smoothed Particle Hydrodynamics), która wykorzystuje podejście znane z dynamiki molekularnej i używa wariantu dynamiki Newtonowskiej z parametrami, które mają symulować zachowanie cieczy realnej. Sprawdzają się jednak w symulacjach w czasie rzeczywistym lub sytuacjach gdzie wymagana jest <<interakcja z płynem>>. Te metody wykorzystywane są również we wspomnianych powyżej gałęziach przemysłu, jednak są przydatne jedynie na etapie projektowania. Poza tym metody online znajdują zastosowanie w symulacjach medycznych lub grach wideo.
\par

\subsection{Mój wkład -- założenia}
\paragraph{}
W niniejszej pracy moim celem było zaprezentowanie metody zdolnej do symulacji płynu w czasie rzeczywistym. Głównym założeniem pracy była możliwość interakcji z płynem w stosunkowo dużej skali przy zachowaniu satysfakcjonującej wydajności. Mój wybór, przy sugestii promotora, padł na metodę, która zdobywa w ostatnim dziesięcioleciu najwięcej rozgłosu: Smoothed Particle Hydrodynamics (pol. wygładzonej hydrodynamiki cząstek).\\
Osobnym krokiem w ramach symulacji płynu jest jego wizualizacja, która sprowadza się do dwóch elementów: wydobywania powierzchni płynu oraz jej renderowania. Ten fragment często pochłania najwięcej czasu symulacji i kluczowym jest znalezienie efektywnego rozwiązania. W niniejszej pracy przedstawiona zostało jedno z możliwych rozwiązań tego problemu.
\par
